\documentclass[border=10pt]{standalone}
\usepackage{ctex}
\usepackage{amsmath,amssymb}
\usepackage{tikz}
\usepackage{tikz-feynman}
\tikzfeynmanset{compat=1.1.0}

\begin{document}
\begin{tikzpicture}

\node[font=\Large\bfseries] at (7, 0) {$\phi^4$ 二阶真空图的三种拓扑};

% ===== 8-shape: 4 internal lines =====
\begin{feynman}
  \vertex (v1a) at (0, -3);
  \vertex (v2a) at (3, -3);
  \diagram* {
    (v1a) -- [bend left=80, edge label=$G_F$] (v2a),
    (v1a) -- [bend left=30] (v2a),
    (v1a) -- [bend right=30] (v2a),
    (v1a) -- [bend right=80, edge label'=$G_F$] (v2a),
  };
  \filldraw[black] (v1a) circle (2.5pt);
  \filldraw[black] (v2a) circle (2.5pt);
\end{feynman}
\node[font=\small] at (1.5, -4.5) {4 条内线};
\node[font=\small, red] at (1.5, -5.0) {对称因子 $1/48$};

% ===== Figure-8: 2 internal + self-loops =====
\begin{feynman}
  \vertex (v1b) at (5.5, -3);
  \vertex (v2b) at (8.5, -3);
  \diagram* {
    (v1b) -- [bend left=30, edge label=$G_F$] (v2b),
    (v1b) -- [bend right=30, edge label'=$G_F$] (v2b),
    (v1b) -- [loop, min distance=15mm, looseness=8, edge label=$G_F(0)$] (v1b),
    (v2b) -- [loop, min distance=15mm, looseness=8, edge label=$G_F(0)$] (v2b),
  };
  \filldraw[black] (v1b) circle (2.5pt);
  \filldraw[black] (v2b) circle (2.5pt);
\end{feynman}
\node[font=\small] at (7, -4.5) {2 条内线 + 自环};
\node[font=\small, red] at (7, -5.0) {对称因子 $1/16$};

% ===== Disconnected: two separate bubbles =====
\begin{feynman}
  \vertex (v1c) at (11, -3);
  \vertex (v2c) at (14, -3);
  \diagram* {
    (v1c) -- [bend left=60, edge label=$G_F$] (v1c),
    (v2c) -- [bend left=60, edge label=$G_F$] (v2c),
  };
  \filldraw[black] (v1c) circle (2.5pt);
  \filldraw[black] (v2c) circle (2.5pt);
\end{feynman}
\node[font=\small] at (12.5, -4.5) {0 条内线（断开图）};
\node[font=\small, red] at (12.5, -5.0) {在 $\ln Z$ 中消失};

% ===== Classification summary =====
\node[font=\small, text width=14cm, align=left] at (7, -6.5) {%
  \textbf{分类}（按 $x_1$-$x_2$ 连线数）：\\[3pt]
  $\bullet$ 4 条连线：两顶点间 4 条传播子 = ``8字形''\\
  $\bullet$ 2 条连线 + 2 自环：两顶点间 2 条 + 各自 tadpole = ``日字形''\\
  $\bullet$ 0 条连线：两个独立真空泡 = 断开图 $\Rightarrow$ 被连通图求和 $\ln Z$ 排除
};

\end{tikzpicture}
\end{document}
